\subsection{Implementácia systému lokálnych notifikácií}

Aplikácia Joinly informuje používateľov o nadchádzajúcich udalostiach
prostredníctvom lokálnych notifikácií systému Android. Na rozdiel od push
notifikácií, ktoré vyžadujú serverovú infraštruktúru, lokálne notifikácie
sú naplánované priamo na zariadení a nevyžadujú aktívne internetové
pripojenie v momente doručenia.

\paragraph{Architektúra systému:}
Systém pozostáva z troch spolupracujúcich komponentov:
\begin{itemize}
    \item \textbf{NotificationHelper} — singleton objekt zodpovedný za
          plánovanie a rušenie notifikácií,
    \item \textbf{NotificationReceiver} — prijímač systémových správ,
          ktorý zostavuje a zobrazuje notifikáciu,
    \item \textbf{Notification} — dátový model zabezpečujúci perzistenciu
          naplánovaných notifikácií vo Firestore.
\end{itemize}

\paragraph{Notifikačný kanál a správa oprávnení:}
Od verzie Android 8.0 (Oreo) je povinné zaregistrovať notifikačný kanál
pred prvým použitím. Kanál s identifikátorom \texttt{joinly\_events}
a prioritou \texttt{IMPORTANCE\_HIGH} sa vytvára pri štarte aplikácie
v \texttt{LaunchedEffect} v kompozícii \texttt{App.kt}:

\begin{lstlisting}[language=Kotlin, caption=Vytvorenie notifikačného kanála]
fun createChannel(context: Context) {
    if (Build.VERSION.SDK_INT >= Build.VERSION_CODES.O) {
        val channel = NotificationChannel(
            CHANNEL_ID, CHANNEL_NAME, NotificationManager.IMPORTANCE_HIGH
        ).apply {
            description = "Upozornenia na nadchadzajuce udalosti"
        }
        context.getSystemService(NotificationManager::class.java)
            .createNotificationChannel(channel)
    }
}
\end{lstlisting}

Od verzie Android 13 (Tiramisu) vyžaduje zobrazovanie notifikácií
explicitné oprávnenie \texttt{POST\_NOTIFICATIONS}. Aplikácia ho vyžiada
pri prvom spustení pomocou \texttt{rememberLauncherForActivityResult}.
Oprávnenie \texttt{SCHEDULE\_EXACT\_ALARM} je deklarované v súbore
\texttt{AndroidManifest.xml} a umožňuje presné načasovanie budíkov.

\paragraph{Plánovanie notifikácie:}
Notifikácia sa naplánuje automaticky v momente, keď sa používateľ pripojí
k udalosti ako potvrdený účastník (nie ako čakajúci v poradovníku).
Upozornenie sa odošle hodinu pred začiatkom udalosti. Plánovanie
zabezpečuje \texttt{AlarmManager} prostredníctvom metódy
\texttt{setExactAndAllowWhileIdle}, ktorá garantuje spustenie aj v prípade,
že je zariadenie v úspornom režime (\textit{Doze mode}). Na zariadeniach
s Android 12+ sa pred použitím presného alarmu overí dostupnosť oprávnenia:

\begin{lstlisting}[language=Kotlin, caption=Plánovanie notifikácie pomocou AlarmManager]
fun scheduleEventNotification(context: Context, event: Event) {
    val dateFrom = event.dateFrom ?: return
    val notifyAt = dateFrom.time - NOTIFY_BEFORE_MS   // 1 hodina pred
    if (notifyAt <= System.currentTimeMillis()) return

    val alarmManager = context.getSystemService(AlarmManager::class.java)
    val pending = buildPendingIntent(context, event)

    if (Build.VERSION.SDK_INT >= Build.VERSION_CODES.S
            && !alarmManager.canScheduleExactAlarms()) {
        alarmManager.setAndAllowWhileIdle(
            AlarmManager.RTC_WAKEUP, notifyAt, pending
        )
    } else {
        alarmManager.setExactAndAllowWhileIdle(
            AlarmManager.RTC_WAKEUP, notifyAt, pending
        )
    }
}
\end{lstlisting}

Každý alarm je identifikovaný pomocou \texttt{PendingIntent}, do ktorého
sú vložené metadáta udalosti — \texttt{eventId}, \texttt{eventTitle},
\texttt{eventPlace} a čas začiatku. Ako kľúč pre \texttt{requestCode}
sa používa \texttt{eventId.hashCode()}, čo zaručuje jednoznačnú
identifikáciu každej naplánovanej notifikácie.

\paragraph{Zobrazenie notifikácie:}
Keď \texttt{AlarmManager} spustí alarm, systém doručí broadcast triede
\texttt{NotificationReceiver}, ktorá implementuje \texttt{BroadcastReceiver}.
Prijímač z \texttt{Intent} extrahuje metadáta udalosti, zostaví text správy
a zobrazí notifikáciu cez \texttt{NotificationManagerCompat}:

\begin{lstlisting}[language=Kotlin, caption=Zostavenie a zobrazenie notifikácie]
override fun onReceive(context: Context, intent: Intent) {
    val eventTitle = intent.getStringExtra("eventTitle") ?: "Udalost"
    val eventPlace = intent.getStringExtra("eventPlace") ?: ""
    val eventTime  = intent.getLongExtra("eventTime", 0L)

    val timeStr = SimpleDateFormat("HH:mm", Locale.getDefault())
        .format(Date(eventTime))
    val body = "\"$eventTitle\" zacina o $timeStr - $eventPlace"

    val notification = NotificationCompat.Builder(
            context, NotificationHelper.CHANNEL_ID
        )
        .setSmallIcon(R.drawable.ic_star)
        .setContentTitle("Joinly -- Event o hodinu!")
        .setContentText(body)
        .setStyle(NotificationCompat.BigTextStyle().bigText(body))
        .setPriority(NotificationCompat.PRIORITY_HIGH)
        .setAutoCancel(true)
        .build()

    NotificationManagerCompat.from(context)
        .notify(eventId.hashCode(), notification)
}
\end{lstlisting}

Trieda \texttt{NotificationReceiver} je registrovaná v súbore
\texttt{AndroidManifest.xml} ako prijímač systémových správ.

\paragraph{Zrušenie notifikácie:}
Ak používateľ opustí udalosť alebo ak organizátor udalosť vymaže,
naplánovaný alarm sa zruší volaním \texttt{cancelEventNotification}.
Zrušenie sa vykoná vyhľadaním existujúceho \texttt{PendingIntent}
s príznakmi \texttt{FLAG\_NO\_CREATE} — ak alarm neexistuje, metóda
nevyvolá chybu:

\begin{lstlisting}[language=Kotlin, caption=Zrušenie naplánovanej notifikácie]
fun cancelEventNotification(context: Context, eventId: String) {
    val alarmManager = context.getSystemService(AlarmManager::class.java)
    val intent = Intent(context, NotificationReceiver::class.java)
    val pending = PendingIntent.getBroadcast(
        context,
        eventId.hashCode(),
        intent,
        PendingIntent.FLAG_NO_CREATE or PendingIntent.FLAG_IMMUTABLE
    )
    pending?.let { alarmManager.cancel(it) }
}
\end{lstlisting}

\paragraph{Perzistencia vo Firestore:}
Okrem lokálneho alarmu sa každá naplánovaná notifikácia uloží aj do
Firestore kolekcie \texttt{notifications}. Dátový model \texttt{Notification}
uchováva \texttt{userId}, \texttt{eventId}, \texttt{eventTitle}, čas
odoslania \texttt{scheduledAt} a príznak \texttt{sent}. Tento záznam sa
pri opustení udalosti alebo vymazaní eventu tiež odstráni. Perzistencia vo
Firestore slúži ako záloha pre prípad, keby sa lokálny alarm stratil
(napríklad po reštarte zariadenia).